\section{Actual points and derivations without degree-zero finiteness}
\label{sec:cubic-model-points-without-m2}

The completion in AC-G56 gives a stronger interface than
the origin calculation alone.  The degree-zero ring of
the actual weighted presentation is complete local, and
the full finite-weight algebra is Noetherian over it.
Those facts give the finite-image and continuity
arguments needed for points and derivatives.
We prove the entire comparison and retain the
distinction between an individual point's finite image
and the global condition \texttt{M2}.

Keep the actual notation
$P=O[[z_1,\ldots,z_d]]$, $R=P/I$, $A$, $J$ and
$S=R/JR$ of AC-G56, where $O$ is the coefficient
discrete valuation ring and all variable weights
belong to $\{-2,-1,0,1,2\}$.
The ring $A$ is the image of all finite-weight series,
not an arbitrarily selected polynomial subalgebra.
Write $A_0$ for its degree-zero ring, and write
$\overline E=\overline{\mathbb Q}_\ell$ with its
specified nonarchimedean absolute value.

\begin{lemma}[AC-G58-algebra: the actual complete zero part and finite images]
\label{AC-G58-algebra}
The ring $A_0$ is a complete Noetherian local
$O$-algebra with residue field $O/\varpi O$.
The ring $A$ is generated by the finitely many
$z_i$ over $A_0$, and is therefore Noetherian.
For every $O$-algebra map
$\lambda:A_0\to\overline E$ there is a finite
extension $L/\operatorname{Frac}(O)$ such that
\[
 \lambda(A_0)\subset\mathcal O_L,\qquad
 \lambda(\mathfrak m_{A_0})\subset\mathfrak m_L,
 \qquad \lambda(A_0)\text{ is a finite }O\text{-module}.
                                                               \tag{G58.1}
\]
The map is automatically continuous.
No condition on $J$ or \texttt{M2} is needed.
\end{lemma}
\begin{proof}
\proofstep{The exact zero part and its algebra generators.}
There are finitely many nonconstant zero-weight
monomials in the following list:
the variables of weight zero; the products of
opposite weights $1,-1$ or $2,-2$; and the products
of weights $(2,-1,-1)$ or $(-2,1,1)$.
Indices may repeat in a triple.
These generate the monoid of all zero-weight
monomials.  Indeed remove weight-zero variables,
then remove opposite equal-weight pairs.
Any remaining mixture has only weight $2$ against
weight $-1$, or only weight $-2$ against weight
$1$; total weight zero then partitions it into
the indicated triples.

Denote the resulting finite list by $M_1,\ldots,M_v$.
Substitution $T_i\mapsto M_i$ gives a map
$O[[T_1,\ldots,T_v]]\to P_0$.
It is well defined because every $M_i$ has positive
total degree.  It is onto: choose one factorization
of each zero-weight monomial into the $M_i$, and
place its coefficient at that chosen exponent
vector in the $T$-series.  Different original
monomials cannot have the same chosen exponent
vector, and each output coefficient receives only
finitely many terms.
The composite onto $A_0$ is consequently
surjective.  Its kernel has zero constant term,
by the actual origin.  Thus $A_0$ is a quotient
of this finite-variable complete Noetherian local
ring and has exactly the asserted local, complete
and residue-field properties.

For a monomial of a fixed weight $b$, remove
zero-weight divisors from the same finite list
until none remain.
There are only finitely many possible remaining
monomials of that weight.
To check this directly, a remaining monomial
either has variables of one sign only, or consists
of weight-$2$ variables and one weight-$(-1)$
variable, or of weight-$(-2)$ variables and one
weight-$1$ variable.
Its total degree is at most $|b|+2$.
Partition the terms of any homogeneous series
according to one such remainder.  It becomes a
finite sum of a polynomial monomial times a
weight-zero series.
After passage to $R$ this proves
$A=A_0[z_1,\ldots,z_d]$ and hence Noetherianity.

These are precisely the literal source facts
\path{CommutativeAlgebra/WPresGrZeroIsCOOfMemMaximalIdeal.lean},
\texttt{wPres\_gr\_zero\_isCO\_of\_mem\_maximalIdeal},
and \path{CommutativeAlgebra/GrGen.lean},
\texttt{gr\_gen}, applied by
\texttt{UpperTheory.V.Alg.isWRing} in
\path{CommutativeAlgebra/IsWRing.lean}.
The explicit monomial proof above verifies their
application for the actual bounded weights.
All paths in this section have the pinned root
\path{Development/AutomorphicArithmetic/AutomorphicVarieties/}.

\proofstep{Automatic locality of algebraic-valued points.}
More generally let $B$ be a local $O$-algebra and
$\lambda:B\to\overline E$ an $O$-algebra map.
For $r\in\mathfrak m_B$, suppose that
$a=\lambda(r)$ has $|a|_\ell\geq1$.
Then $a\ne0$ and $a^{-1}$ is integral over $O$.
Choose its monic equation
$T^n+c_{n-1}T^{n-1}+\cdots+c_0$ over $O$.
After multiplication by $a^n$ it gives
\[
             1+c_{n-1}a+\cdots+c_0a^n=0.
                                                               \tag{G58.2}
\]
But $1+c_{n-1}r+\cdots+c_0r^n$ belongs to
$1+\mathfrak m_B$, so is a unit.  A ring
homomorphism to a field cannot send it to zero.
This contradiction proves
$|\lambda(r)|_\ell<1$.
This automatic-locality fact is already proved
as the private \texttt{pts\_loc} in
\path{UnitaryThreeGroups/ExistsTorusTwistEqLocalPointOfWeightedPres.lean};
the present improvement will use it on the
actual $A_0$.

Now put $B=A_0$.
Every $b\in B$ is congruent modulo
$\mathfrak m_B$ to an element of $O$.
It follows from the just proved inequality that
$|\lambda(b)|_\ell\leq1$.
Choose finitely many generators
$u_1,\ldots,u_s$ of $\mathfrak m_B$, and put
\[
        \rho=\max\bigl(|\varpi|_\ell,
                       |\lambda(u_1)|_\ell,\ldots,
                       |\lambda(u_s)|_\ell\bigr)<1 .
                                                               \tag{G58.3}
\]
All coefficients in an ideal expression have
image of norm at most one.
Hence
$|\lambda(b)|_\ell\leq\rho^N$ for
$b\in\mathfrak m_B^N$, proving continuity.

Let $L$ contain $\operatorname{Frac}(O)$ and
the finitely many $\lambda(u_i)$.
For each $b\in B$ and each $N$, repeated
residue-field approximation gives an
$O$-polynomial $P_N$ in the $u_i$ with
$b-P_N\in\mathfrak m_B^N$.
Explicitly, choose its constant term modulo
$\mathfrak m_B$; express the error by the
$u_i$; approximate the new coefficients
modulo $\mathfrak m_B$ and iterate.
Then $\lambda(P_N)\in L$ and
$\lambda(P_N)\to\lambda(b)$.
A finite extension $L$ is complete and hence
closed in $\overline E$, so
$\lambda(b)\in L$.
It is integral by its norm bound.
Finally $\lambda(B)$ is an $O$-submodule of
the finite $O$-module $\mathcal O_L$.
Since $O$ is Noetherian, it is itself finite.
This proves (G58.1).
It also independently verifies the finite-image
conclusion supplied by the constant-dual-point
instance of \texttt{Loc.ptFin} in
\path{PadicHodgeTheory/PtFinFinite.lean}.
\end{proof}

\begin{theorem}[AC-G58-points: the actual point field without \texttt{M2}]
\label{AC-G58-points}
For every $O$-algebra point $\xi:A\to\overline E$
killing $J$, there are $t\in\overline E^\times$
and an actual local point $f:S\to\overline E$
such that
\[
             \operatorname{act}(t)\xi=f\circ q_A .
                                                               \tag{G58.4}
\]
Here $\operatorname{act}(t)$ multiplies weight
$b$ by $t^b$.
The entire image of $f$ lies in the ring of
integers of one finite coefficient extension.
This proves the actual model's \texttt{pts}
field without its \texttt{M2} antecedent.
\end{theorem}
\begin{proof}
The restriction of $\xi$ to $A_0$ has the
properties of AC-G58-algebra.
Give $\overline E$ the valuation $v$ normalized
by $v(\ell)=1$.
For each coordinate with nonzero image put
$v_i=v(\xi(z_i))\in\mathbb Q$.

If $w_i=a>0$ and $w_j=-b<0$, the monomial
$z_i^b z_j^a$ has weight zero and vanishes
under the origin.
It lies in $\mathfrak m_{A_0}$: a unit there
could not map to zero under the origin.
If both coordinate images are nonzero,
automatic locality therefore gives
\[
                         b v_i+a v_j>0 .
                                                               \tag{G58.5}
\]
The finite set of inequalities
\[
        s>-v_i/w_i\quad(w_i>0),\qquad
        s<v_i/(-w_i)\quad(w_i<0)
                                                               \tag{G58.6}
\]
has a rational solution $s$.
Indeed (G58.5) says that every lower endpoint
is strictly smaller than every upper endpoint.
If either family is empty there is only a
one-sided bound; if both are empty any $s$
works.  Coordinates with zero image impose
no restriction.
Choose $t\in\overline E^\times$ with $v(t)=s$,
which is possible by taking a root of an
integral or negative power of $\ell$.

Set $\xi_t=\operatorname{act}(t)\xi$.
Then $|\xi_t(z_i)|_\ell<1$ for every coordinate
of nonzero weight.
The same holds for a weight-zero coordinate,
because it belongs to $\mathfrak m_{A_0}$.
The map $\xi_t$ agrees with $\xi$ on $A_0$,
and $A$ is generated over $A_0$ by the $z_i$.
Consequently its entire image is contained
in $\mathcal O_L$ for one finite extension
$L$ containing the image field of $A_0$,
the finitely many $\xi(z_i)$, and $t$.

By AC-G56-jets the origin maximal ideal of
$A$ is
$\mathfrak m_A=(\varpi,z_1,\ldots,z_d)$.
Put
$\rho=\max(|\varpi|_\ell,|\xi_t(z_i)|_\ell)<1$.
The image of every coefficient in $A$ has
norm at most one; hence
\[
           |\xi_t(a)|_\ell\leq\rho^N
                       \qquad(a\in\mathfrak m_A^N).
                                                               \tag{G58.7}
\]
Thus $\xi_t$ is continuous and extends uniquely
to the completion $R=\widehat A_{\mathfrak m_A}$,
with values in the complete ring $\mathcal O_L$.
Each generator $g_k$ of $J$ is homogeneous,
so $\xi_t(g_k)=t^{w(g_k)}\xi(g_k)=0$.
The extension therefore kills the full ideal
$JR$ and factors through $S$.
It maps the maximal ideal of $S$, generated
by the images of $\varpi,z_i$, into
$\mathfrak m_L$ and satisfies \texttt{IsLocPt}.
Its restriction is (G58.4).
This proves the actual point comparison.

For the source interface, the old proof in
\path{UnitaryThreeGroups/ExistsTorusTwistEqLocalPointOfWeightedPres.lean}
uses \texttt{M2} through \texttt{fg\_lemma}
only to obtain a finite image of the degree-zero
part of the chosen point.
That finite image is now proved directly by
AC-G58-algebra; the strict rescaling and extension
above make the replacement explicit.
\end{proof}

\begin{lemma}[AC-G58-der: bounded derivatives and the actual dual-point field]
\label{AC-G58-der}
Let $f:S\to\overline E$ be any actual local
point, put $\phi=fq_A$, and let
$D:A\to\overline E$ be an $O$-linear derivation
at $\phi$ which kills $J$.
There is a unique continuous derivation
$D_S:S\to\overline E$ at $f$ restricting to
$D$ along $q_A$.
Its image is contained in a bounded
$O$-module in one finite coefficient field.
Consequently
\[
                 F(s)=f(s)+\varepsilon D_S(s)
                                                               \tag{G58.8}
\]
is an actual dual point on $S$, with
$\operatorname{dualFst}(F)=f$.
Thus the actual model's \texttt{der} field
also holds without \texttt{M2}.
\end{lemma}
\begin{proof}
Since $S$ is a complete Noetherian local
$O$-algebra with the original residue field,
the finite-image argument of AC-G58-algebra,
or the exact constant-dual instance of
\texttt{Loc.ptFin}, applies to $f$.
Thus $f(S)\subset\mathcal O_L$ for a finite
extension $L$, and $\phi(A)$ is a finite
$O$-submodule of $\mathcal O_L$.

The crucial continuity assertion for $D$ is
proved before extending it.
The ring $A$ is Noetherian by AC-G58-algebra,
so $K=\ker\phi$ has finitely many ideal
generators $k_1,\ldots,k_u$.
Choose $b_1,\ldots,b_v\in A$ whose images
generate $\phi(A)$ as an $O$-module.
Every $a\in A$ has an expression
\[
                 a=\sum_i o_i b_i+\sum_j a_jk_j,
                    \qquad o_i\in O,\quad a_j\in A .
                                                               \tag{G58.9}
\]
The derivation rule and $\phi(k_j)=0$ give
\[
                 D(a)=\sum_i o_iD(b_i)
                              +\sum_j\phi(a_j)D(k_j).
                                                               \tag{G58.10}
\]
Adjoin the finitely many $D(b_i),D(k_j)$
to $L$, obtaining a finite extension $L'$.
The right side belongs to the finite
$O$-module generated by the $D(b_i)$ and
by an $O$-basis of $\phi(A)$ times the
$D(k_j)$.
It is bounded in $L'$.
Choose $C\geq1$ bounding $|D(a)|_\ell$
for every $a\in A$.
This is an algebraic boundedness argument,
not an assumption that all derivations of
a dense subring are continuous.

Put
\[
          \rho=\max(|\varpi|_\ell,
                     |\phi(z_1)|_\ell,\ldots,
                     |\phi(z_d)|_\ell)<1 .
                                                               \tag{G58.11}
\]
The strict inequality follows from
\texttt{IsLocPt} and the fact that all
coordinates map to the maximal ideal of
$S$.  For a monomial of degree $N$ in
$\varpi,z_i$, the product rule expresses
its derivative as a sum of terms with
$N-1$ such point values and one derivative.
For an arbitrary coefficient multiple there
is also a term with $N$ point values and
a coefficient derivative.
All coefficient point values have norm at
most one, all coefficient derivatives are
bounded by $C$, and the ultrametric inequality
therefore gives
\[
              |D(a)|_\ell\leq C\rho^{N-1}
                    \qquad(a\in\mathfrak m_A^N,\ N\geq1).
                                                               \tag{G58.12}
\]
This proves adic continuity into the complete
finite field $L'$.

The completion isomorphism of AC-G56-jets
extends $D$ uniquely to $R$.
The point $\phi$ likewise extends, and its
extension is $f$ composed with $R\to S$:
both continuous maps agree on the dense
subring $A$.
Passing to limits in the derivation identity
proves the identity on $R$.
For each generator $g_k$ of $J$, both the
extended point and the extended derivative
vanish at $g_k$.
The derivation rule therefore kills every
element of $JR$, and the extension descends
to $D_S$ on $S$.
It remains in the same finite $O$-module:
such a module is compact, hence closed in
the finite extension $L'$.
Uniqueness among continuous derivations
follows from density of $q_A(A)$.
The derivation rule is precisely the
multiplicativity of (G58.8).
This constructs the actual dual point with
its prescribed first component, rather
than only a point of an abstract cotangent
space.
\end{proof}

\begin{theorem}[AC-G58: \texttt{M3} for the actual cubic model]
\label{AC-G58}
Under the retained hypotheses of the actual
all-good-prime model, \texttt{M3} holds without
assuming \texttt{M2}.
More precisely, the only point of $A/J$
vanishing on every nonzero weight is its
actual origin.
\end{theorem}
\begin{proof}
Let $\xi$ be such a point.
Apply AC-G58-points, obtaining
$\operatorname{act}(t)\xi=fq_A$.
Since $\xi$ is zero in every nonzero weight
and the torus action is the identity in
weight zero, $\operatorname{act}(t)\xi=\xi$.
Every off-diagonal universal entry has a
homogeneous preimage of nonzero weight,
so its value at $f$ is zero.
In particular the actual representation is
both \texttt{Upper} and \texttt{Lower}.
The all-good-prime fixed-point theorem
AC-G55-split gives $f=\overline x_0$.
Thus $\xi=\overline x$.
In the exact definition of \texttt{M3} in
\path{CommutativeAlgebra/OrbitVanishingLocal.lean},
the antecedent asks for a different point
vanishing in both signs.  We have proved
that antecedent impossible.
This is the same fixed-point implication as
\path{UnitaryThreeGroups/GradedModelPointConditionOfPtHyp.lean},
now with the proved \texttt{M2}-free point
comparison.
\end{proof}

\begin{corollary}[AC-G58-M4: the remaining arithmetic rigidity input]
\label{AC-G58-M4}
Assume, in addition, that at every actual
\texttt{TwoPt} of the cubic ring every actual
dual-point tangent cocycle is a coboundary.
This is exactly the still separate
\texttt{PtHyp.rigid} input.
Then the source condition \texttt{M4} holds,
without an \texttt{M2} antecedent.
No such arithmetic rigidity is asserted by
this corollary.
\end{corollary}
\begin{proof}
Take a point $\xi$ of $A/J$ nonzero in both
weight signs, and two derivations $D,D'$ at
$\xi$ killing $J$.
Use AC-G58-points to replace $\xi$ by
$\xi_t=fq_A$.  On a homogeneous element of
weight $b$ define
$D_t(a)=t^bD(a)$ and $D'_t(a)=t^bD'(a)$,
and extend by the finite weight decomposition.
The product rule and $t^{b+c}=t^bt^c$
show these are derivations at $\xi_t$.
They kill $J$, since its generators are
homogeneous and vanish at $\xi_t$.
Lemma AC-G58-der gives actual dual points
$F,F'$ above $f$ with these derivatives.

The point $f$ is an actual \texttt{TwoPt}.
For example, if its representation were upper
triangular, all negative-weight entries would
vanish, on the index-two subgroup and hence
on the full semidirect group because the
matrix at conjugation is $I$.
Every negative-weight entry monomial would
then vanish.  The homogeneous entry density
of AC-G51 and continuity of $f$ would make
$fq_A$ zero on all negative weights,
contrary to $\xi$ being nonzero there.
The other triangular case is identical.

By the additional rigidity hypothesis,
$z_F=\partial X$ and $z_{F'}=\partial X'$.
The actual normalization makes these
normalized cocycles.
Its marked parameter is a unit: modulo
the maximal ideal its \texttt{zeta} value is
the reduction of
\[
          \chi(g_1)^j
                    \bigl(\gamma(g_1)-\gamma(g_1)^{-1}\bigr),
                                                               \tag{G58.13}
\]
which is a unit by the retained mark.
The full seed-independent
\texttt{normCoc\_cob\_iff} in
\path{UnitaryThreeGroups/NormCob.lean}
therefore puts both $X$ and $X'$ in the
one-dimensional torus Lie algebra
$\overline E\operatorname{diag}(1,0,-1)$.
Thus some $(a,a')\ne(0,0)$ satisfies
$a z_F+a'z_{F'}=0$.
Multiplying the cocycle identity by the
common specialized representation matrix
shows that the same linear combination
of derivatives vanishes on every actual
universal entry.
It vanishes on their dense $O$-algebra and
then, by the continuity proved above, on
$S$.
Restricting to $A$ gives
$aD_t+a'D'_t=0$.
On each homogeneous part $t^b\ne0$, so
$aD+a'D'=0$.
This is exactly \texttt{M4}.
It also rechecks the two extension steps in
the full proof
\path{UnitaryThreeGroups/GradedModelDerivationConditionOfPtHyp.lean}
with their \texttt{M2} antecedents removed
by the preceding results.
\end{proof}

\paragraph{What remains global.}
The conclusion (G58.1) concerns the image
of each individual point and allows its
finite coefficient field and generators to
depend on that point.
It does not prove that the whole degree-zero
quotient $\operatorname{im}(A_0\to A/J)$
is finite over $O$, which is the source's
\texttt{M2}.
That condition remains in the numerical
\texttt{LemmaM}; no replacement for it, no
uniform point count and no Selmer-length
upper bound is inferred here.
The new unconditional conclusion is
\texttt{M3}, alongside the already proved
\texttt{M1}; the arithmetic \texttt{M4}
input remains actual nonbase-point rigidity.

\paragraph{Geometric dictionary.}
The argument compares algebraic-valued points
of the finite-weight algebra with points and
dual points of its actual formal completion.
The torus rescales the prescribed matrix
weights, and the derivation estimate justifies
passage to the completion.
The finite local systems being deformed
retain the already audited SGA~1,
Expos\'e~V, \S\S4 and~7 interpretation.
These explicit algebraic and topological
comparisons supply the present conclusions;
the finite-cover dictionary alone does not
provide degree-zero finiteness or arithmetic
rigidity.
